\documentclass[11pt,a4paper,fontset=fandol]{ctexart}

\usepackage[margin=22mm]{geometry}
\usepackage{amsmath,amssymb,bm}
\usepackage{array,longtable}
\usepackage{hyperref}
\usepackage{xcolor}
\usepackage{seqsplit}

\hypersetup{
  colorlinks=true,
  linkcolor=blue!55!black,
  urlcolor=blue!55!black,
  citecolor=blue!55!black
}

\newcommand{\dd}{\mathrm{d}}
\newcommand{\ve}{ve}
\newcommand{\tr}{tr}
\newcommand{\kk}{\bm{k}}
\newcommand{\nn}{\bm{n}}
\newcommand{\uu}{\bm{u}}
\newcommand{\zz}{\bm{z}}
\newcommand{\QQ}{\bm{Q}}
\newcommand{\FF}{\bm{F}}
\newcommand{\RR}{\bm{R}}
\newcommand{\LL}{\bm{L}}
\newcommand{\Gammaone}{\Gamma}
\newcommand{\Tmode}{\mathcal{T}}
\newcommand{\Mve}{\mathcal{M}_{\ve}}
\newcommand{\code}[1]{\texttt{\seqsplit{#1}}}

\title{COREFL 双温无粘通量雅可比、特征 WENO 与源码审查}
\author{COREFL 双温实现审查文档}
\date{2026-05-11}

\begin{document}
\sloppy
\maketitle

\section{审查范围与参考}

本文审查 COREFL 双温模式下的有限差分无粘通量、特征 WENO 投影、振动--电子能量方程、源项和粘性项。
对照来源包括 SU2 v7 NEMO thermochemical nonequilibrium 文档、仓库内
\texttt{化学反应流动无粘通量雅可比矩阵及其特征向量.pdf}，
以及当前源码。
仓库内 \texttt{无粘通量离散.pdf} 目前无法被 \texttt{pdfinfo} 正常解析，
其 WENO 部分不纳入公式核对依据；具体算法以可读的雅可比 PDF、
SU2 方程和当前 \texttt{src/WENO.cu} 为准。

SU2 v7 的 NEMO 文档给出的双温热力学核心为
\[
\rho e^{\ve}=\sum_s \rho_s\left(e_s^{vib}+e_s^{el}\right),
\]
并给出 Landau--Teller/Park 形式的振动--电子松弛源
\[
\dot\Theta^{\tr:\ve}
=\sum_s \rho_s \frac{e_s^{\ve *}-e_s^{\ve}}{\tau_s}.
\]
同页还写出物种扩散通量、各能量模态 Fourier 热通量
\[
\bm{q}^{k}=\kappa^k \nabla T^k,
\]
以及 Eucken 形式中
\(\kappa_s^{\tr}\) 和 \(\kappa_s^{\ve}\) 应分开计算。
这些式子是本文粘性项和源项审查的准则。

\section{守恒变量与双温能量拆分}

COREFL 的双温版本通过编译期开关
\code{COREFL\_ENABLE\_TWO\_TEMPERATURE} 选择，并在设备代码中用
\texttt{if constexpr (kTwoTemperature)} 裁剪路径。
有限差分通量使用的守恒变量写为
\[
\QQ =
\begin{bmatrix}
\rho & \rho u & \rho v & \rho w & \rho E &
\rho\phi_1 & \cdots & \rho\phi_M
\end{bmatrix}^{T}.
\]
其中 transported scalars 为
\[
\phi_a \in \{Y_1,\ldots,Y_{N_s},E_{\ve},\ldots\}.
\]
在双温有限率化学路径中，
\[
i_{\ve}=N_s,\qquad i_{\ve}^{cv}=5+N_s,
\]
即 \texttt{Eve} 是物种之后的第一个 transported scalar，
守恒量为
\[
Q_{5+i_{\ve}}=\rho E_{\ve}.
\]
\texttt{Tve} 只保存在辅助数组和输出变量中，不是独立守恒变量。
本文后续把 \(E_{\ve}\) 始终视为守恒变量中的独立比能标量：
\[
E_{\ve}=\frac{q_{\ve}}{\rho},\qquad q_{\ve}=Q_{5+i_{\ve}}.
\]
它不写作 \(T\) 或 \(T_{\ve}\) 的函数。\(T_{\ve}\) 是为了计算
化学源、VT 松弛和粘性扩散而从 \(E_{\ve}\) 与组分反解出的辅助温度。

\subsection{\texorpdfstring{模态能量映射与 \(E_{\ve}\leftrightarrow T_{\ve}\)}{模态能量映射与 Eve-Tve}}

代码中的物种振动--电子模态能量求值函数记为
\(\varepsilon_{\ve,s}(\Tmode)\)。这里的 \(\Tmode\) 是求值温度，不是
守恒变量。它定义为
\[
\varepsilon_{\ve,s}(\Tmode)=\varepsilon_{vib,s}(\Tmode)+\varepsilon_{el,s}(\Tmode).
\]
若物种给出了特征振动温度 \(\theta_{v,s}>0\)，则
\[
\varepsilon_{vib,s}(\Tmode)=
R_s\frac{\theta_{v,s}}{\exp(\theta_{v,s}/\Tmode)-1},
\]
\[
c_{v,vib,s}(\Tmode)=
R_s
\frac{(\theta_{v,s}/\Tmode)^2\exp(\theta_{v,s}/\Tmode)}
{\left[\exp(\theta_{v,s}/\Tmode)-1\right]^2}.
\]
若 \(\theta_{v,s}\le0\)，代码令该物种振动能为零。

电子能级按配分函数计算。设电子能级温度和简并度为
\(\theta_{el,sj},g_{sj}\)，
\[
Z_{el,s}(\Tmode)=\sum_j g_{sj}\exp(-\theta_{el,sj}/\Tmode),
\]
\[
\varepsilon_{el,s}(\Tmode)=
R_s\frac{\sum_jg_{sj}\theta_{el,sj}
\exp(-\theta_{el,sj}/\Tmode)}
{Z_{el,s}(\Tmode)},
\]
\[
c_{v,el,s}(\Tmode)=
\frac{R_s}{\Tmode^2}
\left[
\frac{\sum_jg_{sj}\theta_{el,sj}^2
\exp(-\theta_{el,sj}/\Tmode)}{Z_{el,s}(\Tmode)}
-\left(
\frac{\sum_jg_{sj}\theta_{el,sj}
\exp(-\theta_{el,sj}/\Tmode)}{Z_{el,s}(\Tmode)}
\right)^2
\right],
\]
因此
\[
c_{v,\ve,s}(\Tmode)=c_{v,vib,s}(\Tmode)+c_{v,el,s}(\Tmode).
\]

定义混合物模态能量映射
\[
\Mve(\Theta,Y)=\sum_sY_s\varepsilon_{\ve,s}(\Theta).
\]
辅助温度 \(T_{\ve}\) 由隐式关系
\[
\boxed{\Mve(T_{\ve},Y)=E_{\ve}}.
\]
确定。给定守恒标量 \(E_{\ve}\) 和组分 \(Y\)，代码通过 Newton 迭代求
\[
T_{\ve}=\Mve^{-1}(E_{\ve};Y).
\]
反过来，如果边界、初场或 restart 只给出 \(T_{\ve}\)，代码先用同一个映射
\[
E_{\ve}\leftarrow \Mve(T_{\ve},Y)
\]
初始化守恒标量。该赋值属于初始化/同步步骤；在无粘特征推导中
\(E_{\ve}\) 是独立 transported scalar，压力微分不对 \(T_{\ve}\) 求导。

\subsection{NASA 单温热力学拆分}

NASA 单温热力学给出的平衡内能为
\[
e_{eq,s}(T)=h_{eq,s}(T)-R_sT .
\]
它包含该 NASA 拟合所代表的平衡振动--电子贡献。为了在双温状态退化为
热平衡、即 \(T_{\ve}=T\) 且守恒 VE 标量取混合物平衡模态能量值时恢复
NASA 单温热力学，必须从
\(e_{eq,s}(T)\) 中减去同一个平衡温度 \(T\) 下的
\(\varepsilon_{\ve,s}(T)\)，定义平动--转动部分
\[
e_{\tr,s}(T)=e_{eq,s}(T)-\varepsilon_{\ve,s}(T)
=h_{eq,s}(T)-R_sT-\varepsilon_{\ve,s}(T).
\]
式中的 \(\varepsilon_{\ve,s}(T)\) 是 NASA 单温热力学拆分时在平衡温度
\(T\) 下的模态能量求值；它不是守恒变量 \(E_{\ve}\)。双温物种内能可写为
\[
e_s=e_{\tr,s}(T)+\varepsilon_{\ve,s}(T_{\ve}),
\]
但混合物总能量闭合在守恒变量层面写成
\[
E=K+\sum_sY_s e_{\tr,s}(T)+E_{\ve},
\qquad
K=\frac{1}{2}(u^2+v^2+w^2).
\]
这对应 \texttt{src/Thermo.cu} 中
\code{compute\_mixture\_tr\_energy} 和
\texttt{src/FieldOperation.cuh} 中的总能量构造：源码先从
\(h_{eq,s}(T)-R_sT\) 中减去 \(\varepsilon_{\ve,s}(T)\)，再加上独立守恒标量
\texttt{Eve}。

\section{曲线坐标无粘通量}

令度量方向为
\[
\kk=(k_x,k_y,k_z)^T,\qquad |\kk|=\sqrt{k_x^2+k_y^2+k_z^2},
\qquad \nn=\frac{\kk}{|\kk|},
\]
反变速度为
\[
U_k=\kk\cdot\uu=k_xu+k_yv+k_zw .
\]
无粘通量为
\[
\FF_k(\QQ)=
\begin{bmatrix}
\rho U_k\\
\rho uU_k+k_xp\\
\rho vU_k+k_yp\\
\rho wU_k+k_zp\\
(\rho E+p)U_k\\
\rho\phi_1U_k\\
\vdots\\
\rho\phi_MU_k
\end{bmatrix}.
\]
因此双温标量无粘通量严格为
\[
F_{\rho E_{\ve}}^{inv}=\rho E_{\ve}U_k .
\]
在 \texttt{src/WENO.cu} 和 Roe/HLLC/AUSM 路径中，标量通量均按
\(\rho\phi_a U_k\) 处理；当 \(a=i_{\ve}\) 时即为上式。

\section{压力微分和热力学系数}

特征投影需要 \(\partial p/\partial \QQ\)。本节把物种系数和双温系数从
状态方程逐项推出。记
\[
q_s=\rho Y_s,\qquad q_{\ve}=\rho E_{\ve}.
\]
平动--转动内能密度为
\[
\rho e_{\tr}=\rho E-\frac{(\rho u)^2+(\rho v)^2+(\rho w)^2}{2\rho}-q_{\ve}.
\]
因此
\[
\dd(\rho e_{\tr})
=\dd(\rho E)-\uu\cdot\dd(\rho\uu)+K\dd\rho-\dd q_{\ve}.
\]
另一方面，
\[
\rho e_{\tr}=\sum_s q_s e_{\tr,s}(T),
\qquad
p=T\sum_s q_sR_s .
\]
这里把 \(q_s=\rho Y_s\) 作为守恒变量中的物种密度来求微分，
\(T\) 作为由能量闭合反求的热力学变量。对第一式直接使用乘积微分：
\[
\begin{aligned}
\dd(\rho e_{\tr})
&=\dd\left[\sum_s q_s e_{\tr,s}(T)\right] \\
&=\sum_s e_{\tr,s}(T)\,\dd q_s
  +\sum_s q_s \frac{\partial e_{\tr,s}}{\partial T}\,\dd T .
\end{aligned}
\]
由
\[
e_{\tr,s}(T)=h_{eq,s}(T)-R_sT-\varepsilon_{\ve,s}(T)
\]
可得物种平动--转动定容比热
\[
c_{v,\tr,s}(T)
=\frac{\partial e_{\tr,s}}{\partial T}
=c_{p,eq,s}(T)-R_s-c_{v,\ve,s}(T).
\]
混合物平动--转动定容比热定义为
\[
c_{v,\tr}=\sum_sY_s c_{v,\tr,s},
\]
因此
\[
\sum_s q_s \frac{\partial e_{\tr,s}}{\partial T}
=\rho\sum_sY_s c_{v,\tr,s}
=\rho c_{v,\tr}.
\]
于是
\[
\dd(\rho e_{\tr})
=\rho c_{v,\tr}\dd T+\sum_s e_{\tr,s}\dd q_s,
\]
\[
\dd p=\rho R_{mix}\dd T+T\sum_sR_s\dd q_s,
\qquad
R_{mix}=\sum_sY_sR_s .
\]
用
\[
\Gammaone=\gamma_{\tr}-1=\frac{R_{mix}}{c_{v,\tr}},
\qquad
c^2=\gamma_{\tr}R_{mix}T
\]
消去 \(\dd T\) 后即可得到压力对守恒变量的微分。WENO 特征矩阵要求
\(\dd p\) 写成守恒变量
\[
\dd\rho,\quad \dd(\rho u),\quad \dd(\rho v),\quad
\dd(\rho w),\quad \dd(\rho E),\quad \dd q_s,\quad \dd q_{\ve}
\]
的线性组合，因此以下推导先把温度微分完全消去。
首先从守恒变量精确计算 \(\dd(\rho e_{\tr})\)。令
\[
\bm{m}=\rho\uu,\qquad
\rho K=\frac{\bm{m}\cdot\bm{m}}{2\rho}.
\]
动能密度的微分为
\[
\begin{aligned}
\dd(\rho K)
&=\dd\left(\frac{\bm{m}\cdot\bm{m}}{2\rho}\right)\\
&=\frac{\bm{m}}{\rho}\cdot\dd\bm{m}
  -\frac{\bm{m}\cdot\bm{m}}{2\rho^2}\dd\rho\\
&=\uu\cdot\dd\bm{m}-K\dd\rho
 =\uu\cdot\dd(\rho\uu)-K\dd\rho .
\end{aligned}
\]
又因为
\[
\rho e_{\tr}=\rho E-\rho K-q_{\ve},
\]
所以
\[
B=\dd(\rho E)-\uu\cdot\dd\bm{m}+K\dd\rho-\dd q_{\ve}
\]
实际上就是
\[
\boxed{B=\dd(\rho e_{\tr})}.
\]
另一方面，热力学表达已经给出
\[
B=\rho c_{v,\tr}\dd T+\sum_s e_{\tr,s}\dd q_s .
\]
解出 \(\dd T\)：
\[
B-\sum_s e_{\tr,s}\dd q_s
=\rho c_{v,\tr}\dd T,
\]
\[
\boxed{
\dd T
=\frac{B}{\rho c_{v,\tr}}
-\sum_s\frac{e_{\tr,s}}{\rho c_{v,\tr}}\dd q_s
}
\]
或者等价地写成
\[
\dd T
=\frac{1}{\rho c_{v,\tr}}\left(B-\sum_s e_{\tr,s}\dd q_s\right).
\]
接下来逐项代入状态方程微分
\[
\dd p=\rho R_{mix}\dd T+T\sum_sR_s\dd q_s,
\]
得到
\[
\begin{aligned}
\dd p
&=\rho R_{mix}
\frac{1}{\rho c_{v,\tr}}
\left(B-\sum_s e_{\tr,s}\dd q_s\right)
+T\sum_sR_s\dd q_s\\
&=\frac{R_{mix}}{c_{v,\tr}}B
-\frac{R_{mix}}{c_{v,\tr}}\sum_s e_{\tr,s}\dd q_s
+T\sum_sR_s\dd q_s\\
&=\Gammaone B
 -\Gammaone\sum_s e_{\tr,s}\dd q_s
 +\sum_sR_sT\dd q_s\\
&=\Gammaone B
+\sum_s\left(R_sT-\Gammaone e_{\tr,s}\right)\dd q_s .
\end{aligned}
\]
其中最后一步使用
\[
-\Gammaone\sum_s e_{\tr,s}\dd q_s
+\sum_sR_sT\dd q_s
=\sum_s\left(R_sT-\Gammaone e_{\tr,s}\right)\dd q_s .
\]
最后把 \(B\) 展开，得到对守恒变量微分的完整形式：
\[
\boxed{
\dd p
=\Gammaone\left[\dd(\rho E)-\uu\cdot\dd(\rho\uu)+K\dd\rho-\dd q_{\ve}\right]
+\sum_s\left(R_sT-\Gammaone e_{\tr,s}\right)\dd q_s .
}
\]
把每一类守恒变量的系数逐项写出，得到
\[
\begin{aligned}
\dd p
&=\Gammaone K\,\dd\rho
  -\Gammaone u\,\dd(\rho u)
  -\Gammaone v\,\dd(\rho v)
  -\Gammaone w\,\dd(\rho w)
  +\Gammaone\,\dd(\rho E)\\
&\quad
  +\sum_s\left(R_sT-\Gammaone e_{\tr,s}\right)\dd q_s
  -\Gammaone\,\dd q_{\ve}.
\end{aligned}
\]
压力微分没有 \(\dd T_{\ve}\) 项，因为 \(q_{\ve}=\rho E_{\ve}\)
是独立守恒量；\(T_{\ve}\) 只由 \(E_{\ve}\) 和组分反解，用于源项、
粘性项和输出。

COREFL 对所有 transported scalar 使用同一段 WENO 投影循环。为使公式下标
和源码数组下标一致，将物种、\(\rho E_{\ve}\) 以及其他 transported scalar
的压力系数统一记为 \(\alpha_a\)。标量守恒量按
\[
q_a=\rho\phi_a,\qquad
\phi_a\in\{Y_1,\ldots,Y_{N_s},E_{\ve},\text{其他 transported scalar}\}
\]
存储和重构，WENO 投影代码中的 \code{Q[5+a]} 正是 \(q_a\)。
对任意标量 \(q_a\)，定义
\[
\boxed{c^2\alpha_a=
\text{\(\dd q_a\) 在 \(\dd p\) 中的系数}} .
\]
对物种标量，上一式已经给出 \(\dd q_s\) 的系数；
对双温标量，\(\dd q_{\ve}\) 的系数来自
\(-\Gammaone\dd q_{\ve}\)；对不进入压力闭合的其他标量，压力系数为零。
于是所有标量贡献合并为
\[
c^2\sum_a\alpha_a\dd q_a
\]
并由同一个数组 \code{scalar\_alpha[a]} 进入投影。
于是压力微分可写为
\[
\dd p
=\Gammaone\left[\dd(\rho E)-\uu\cdot\dd(\rho\uu)+K\dd\rho\right]
+c^2\sum_a\alpha_a\dd q_a.
\]
此时 \(-\Gammaone\dd q_{\ve}\) 已由求和中的
\(a=\ve\) 项表示，故括号中不再显式写该项。
逐项对应关系为
\[
\begin{array}{c|c|c}
\text{标量守恒量 }q_a & \text{\(\dd p\) 中乘在 \(\dd q_a\) 前的系数}
& \alpha_a\\ \hline
q_s=\rho Y_s
& R_sT-\Gammaone e_{\tr,s}
& (R_sT-\Gammaone e_{\tr,s})/c^2\\
q_{\ve}=\rho E_{\ve}
& -\Gammaone
& -\Gammaone/c^2\\
\text{不参与压力闭合的其他标量}
& 0
& 0
\end{array}
\]
该记号只把
\(\sum_s(\cdots)\dd q_s-\Gammaone\dd q_{\ve}\)
整理为按标量下标 \(a\) 遍历的数组形式；物理变量定义不变：
\(E_{\ve}\) 是独立标量，\(T_{\ve}\) 不进入状态方程。
后面的特征投影统一使用 \(\sum_a\alpha_a\delta(\rho\phi_a)\)。

物种 \(Y_s\) 的系数来自
\[
R_sT-\Gammaone e_{\tr,s}
=\gamma_{\tr}R_sT-\Gammaone h_{\tr,s},
\qquad h_{\tr,s}=e_{\tr,s}+R_sT,
\]
等号右边的推导为
\[
\begin{aligned}
R_sT-\Gammaone e_{\tr,s}
&=R_sT-\Gammaone(h_{\tr,s}-R_sT)\\
&=(1+\Gammaone)R_sT-\Gammaone h_{\tr,s}\\
&=\gamma_{\tr}R_sT-\Gammaone h_{\tr,s}.
\end{aligned}
\]
因此
\[
\boxed{\alpha_s=
\frac{\gamma_{\tr}R_sT-\Gammaone h_{\tr,s}}{c^2}}.
\]
双温标量 \(E_{\ve}\) 只通过
\(-q_{\ve}\) 进入 \(\rho e_{\tr}\)，没有 \(R_sT\) 型状态方程项，因此
\[
c^2\alpha_{\ve}=-\Gammaone,
\]
也就是
\[
\boxed{\alpha_{\ve}=-\frac{\Gammaone}{c^2}}.
\]
不参与压力闭合的被动标量或湍流标量有 \(\alpha_a=0\)。

由这些压力系数定义
\[
\boxed{
\chi_a=-\frac{\alpha_a c^2}{\Gammaone}
}
\]
是该标量波的总能量补偿系数。对物种标量，
\[
\begin{aligned}
\chi_s
&=-\frac{1}{\Gammaone}
  \left(\gamma_{\tr}R_sT-\Gammaone h_{\tr,s}\right)\\
&=h_{\tr,s}-\frac{\gamma_{\tr}R_sT}{\Gammaone}.
\end{aligned}
\]
对双温标量，
\[
\chi_{\ve}
=-\frac{1}{\Gammaone}\left(-\Gammaone\right)=1.
\]

实际守恒变量是 \(q_{\ve}=\rho E_{\ve}\)。推导中的
\(e_{\tr,s}(T)\) 只来自 NASA 单温热力学的平动--转动拆分。
总能量闭合若写成
\(K+\sum_sY_s[h_{eq,s}(T)-R_sT-\varepsilon_{\ve,s}(T_{\ve})]+E_{\ve}\)，
该写法会把独立 VE 能量先减后加，改变
\(\rho e_{\tr}=\rho E-\rho K-\rho E_{\ve}\) 以及压力微分。

以上系数对应 \texttt{src/Thermo.cuh} 中的特征热力学辅助函数。
该函数令 \texttt{scalar\_alpha[s]} 存储 \(\alpha_s\)，
\texttt{scalar\_alpha[i\_eve]} 存储 \(-\Gammaone/c^2\)，
\texttt{energy\_coeff[a]} 存储 \(\chi_a\)。

\[
\begin{array}{c|c|c}
\hline
\text{源码变量} & \text{数学量} & \text{双温标量取值}\\
\hline
\code{scalar\_alpha[l]} & \alpha_l & -\Gammaone/c^2\\
\code{energy\_coeff[l]} & \chi_l=-\alpha_lc^2/\Gammaone & 1\\
\code{Q[5+l]} & q_l=\rho\phi_l & \rho E_{\ve}\\
\code{svm[l]} & \bar\phi_l & \bar E_{\ve}\\
\code{fChar[5+l]} & f_{\phi_l} & f_{E_{\ve}}\\
\hline
\end{array}
\qquad (l=i_{\ve})
\]

\section{线性化特征问题}

先在单位法向 \(\nn\) 上推导，最后再乘以 \(|\kk|\)。
法向通量 \(\FF_n\) 满足 \(\FF_k=|\kk|\FF_n\)，
法向速度为 \(U=\nn\cdot\uu\)。
守恒变量的单位法向无粘通量为
\[
\FF_n=
\begin{bmatrix}
\rho U\\
\rho uU+n_xp\\
\rho vU+n_yp\\
\rho wU+n_zp\\
(\rho E+p)U\\
\rho\phi_1U\\
\vdots\\
\rho\phi_MU
\end{bmatrix}.
\]
对任意扰动，通量微分为
\[
\delta F_{\rho}=\delta\bm{m}\cdot\nn,
\]
\[
\delta\bm{F}_{\rho\uu}
=U\,\delta\bm{m}
+\uu(\delta\bm{m}\cdot\nn-U\delta\rho)+\nn\,\delta p,
\]
\[
\delta F_{\rho E}
=U\,\delta(\rho E)+H(\delta\bm{m}\cdot\nn-U\delta\rho)+U\,\delta p,
\qquad H=E+\frac{p}{\rho},
\]
\[
\delta F_{\rho\phi_a}
=U\,\delta(\rho\phi_a)+\phi_a(\delta\bm{m}\cdot\nn-U\delta\rho).
\]
将特征方程写成
\((\partial\FF_n/\partial\QQ)\delta\QQ=\lambda_n\delta\QQ\)，
并从每一行减去 \(U\delta\QQ\)，即可得到
\(\mu=\lambda_n-U\) 形式。
设扰动为
\[
\delta\QQ=
\left[\delta\rho,\delta\bm{m},\delta(\rho E),
\delta(\rho\phi_1),\ldots,\delta(\rho\phi_M)\right]^T,
\qquad \bm{m}=\rho\uu .
\]
令 \(\mu=\lambda_n-U\)。由
\((\partial\FF_n/\partial\QQ-U I)\delta\QQ=\mu\delta\QQ\)
可得五类基本关系：
\[
\delta\bm{m}\cdot\nn-U\delta\rho=\mu\,\delta\rho,
\]
\[
\uu(\delta\bm{m}\cdot\nn-U\delta\rho)+\nn\,\delta p=\mu\,\delta\bm{m},
\]
\[
H(\delta\bm{m}\cdot\nn-U\delta\rho)+U\,\delta p=\mu\,\delta(\rho E),
\]
\[
\phi_a(\delta\bm{m}\cdot\nn-U\delta\rho)=\mu\,\delta(\rho\phi_a),
\]
\[
\delta p
=\Gammaone\left[
\delta(\rho E)-\uu\cdot\delta\bm{m}+K\delta\rho\right]
+c^2\sum_a\alpha_a\delta(\rho\phi_a).
\]

\subsection{声波}

若 \(\mu\ne0\)，由质量和标量方程得
\[
\delta\bm{m}\cdot\nn-U\delta\rho=\mu\delta\rho,
\qquad
\delta(\rho\phi_a)=\phi_a\delta\rho .
\]
取归一化 \(\delta\rho=1\)。动量方程给出
\[
\delta\bm{m}=\uu+\frac{\delta p}{\mu}\nn .
\]
能量方程给出
\[
\delta(\rho E)=H+\frac{U\,\delta p}{\mu}.
\]
代入压力微分，并利用
\[
c^2\sum_a\alpha_a\phi_a
=c^2-\Gammaone\left(H-K\right),
\]
该恒等式可由热力学定义直接验证。由
\[
\sum_a\alpha_a\phi_a
=\sum_s\alpha_sY_s+\alpha_{\ve}E_{\ve}
\]
和上节给出的系数，有
\[
\begin{aligned}
c^2\sum_a\alpha_a\phi_a
&=\sum_sY_s(\gamma_{\tr}R_sT-\Gammaone h_{\tr,s})
  -\Gammaone E_{\ve}\\
&=\gamma_{\tr}R_{mix}T
  -\Gammaone\left(\sum_sY_sh_{\tr,s}+E_{\ve}\right).
\end{aligned}
\]
又
\[
c^2=\gamma_{\tr}R_{mix}T,\qquad
H-K=\sum_sY_sh_{\tr,s}+E_{\ve},
\]
因此得到上述恒等式。将声波关系
\[
\delta\bm{m}=\uu+\frac{\delta p}{\mu}\nn,\qquad
\delta(\rho E)=H+\frac{U\delta p}{\mu},\qquad
\delta(\rho\phi_a)=\phi_a
\]
代入压力微分，括号内的动量和能量压力项相互抵消：
\[
\begin{aligned}
\delta p
&=\Gammaone\left[
H+\frac{U\delta p}{\mu}
-\uu\cdot\left(\uu+\frac{\delta p}{\mu}\nn\right)+K
\right]
+c^2\sum_a\alpha_a\phi_a\\
&=\Gammaone(H-K)+c^2\sum_a\alpha_a\phi_a=c^2 .
\end{aligned}
\]
得到
\[
\delta p=c^2 .
\]
再由质量关系和动量关系的法向分量可得
\[
\nn\cdot\delta\bm{m}-U\delta\rho=\mu\delta\rho,
\]
\[
\nn\cdot\delta\bm{m}
=U\delta\rho+\frac{\delta p}{\mu},
\]
两式相减得到
\[
\mu^2=\delta p=c^2.
\]
因此单位法向声波特征值为
\[
\lambda_n=U-c,\qquad \lambda_n=U+c.
\]
对应右特征向量为
\[
r_{\pm}=
\begin{bmatrix}
1\\
u\pm cn_x\\
v\pm cn_y\\
w\pm cn_z\\
H\pm cU\\
\phi_1\\
\vdots\\
\phi_M
\end{bmatrix}.
\]

\subsection{切向波、接触波和标量波}

若 \(\mu=0\)，质量方程给出 \(\delta U=0\)，动量方程给出
\(\delta p=0\)。因此 \(\lambda_n=U\) 至少有切向、接触和标量波。
这里的 \(\delta U\) 表示速度扰动的法向分量
\[
\delta U=\nn\cdot(\delta\bm{m}/\rho-\uu\,\delta\rho/\rho).
\]
在下列三个族中均有
\[
\delta\bm{m}\cdot\nn-U\delta\rho=0,
\qquad \delta p=0,
\]
因此满足 \(\mu=0\) 的质量、动量法向、能量和标量方程。

取两组单位切向量 \(\bm{t},\bm{b}\)，满足
\[
\nn\cdot\bm{t}=\nn\cdot\bm{b}=\bm{t}\cdot\bm{b}=0,
\qquad
\bm{b}=\nn\times\bm{t}.
\]
记 \(U_t=\bm{t}\cdot\uu\), \(U_b=\bm{b}\cdot\uu\)。切向波为
\[
r_t=
\begin{bmatrix}
0\\t_x\\t_y\\t_z\\U_t\\0\\\vdots\\0
\end{bmatrix},
\qquad
r_b=
\begin{bmatrix}
0\\b_x\\b_y\\b_z\\U_b\\0\\\vdots\\0
\end{bmatrix}.
\]

接触波取
\[
\delta\rho=1,\qquad \delta\bm{m}=\uu,\qquad
\delta(\rho\phi_a)=\phi_a .
\]
\(\delta p=0\) 要求
\[
0=\Gammaone\left[\delta(\rho E)-K\right]
+c^2\sum_a\alpha_a\phi_a ,
\]
所以接触波能量分量为
\[
\delta(\rho E)
=K-\frac{c^2}{\Gammaone}\sum_a\alpha_a\phi_a
=K+\sum_a\phi_a\chi_a
=H-\frac{c^2}{\Gammaone},
\]
其中
\[
\boxed{\chi_a=-\frac{\alpha_a c^2}{\Gammaone}}.
\]
因此
\[
r_c=
\begin{bmatrix}
1\\u\\v\\w\\H-\frac{c^2}{\Gammaone}\\\phi_1\\\vdots\\\phi_M
\end{bmatrix}.
\]
只有单组分理想气体或没有热力学标量时，
\(H-c^2/\Gammaone\) 才退化为 \(K\)；对反应混合气，
该分量不是 \(K-c^2/\Gammaone\)。

每个 transported scalar 的独立标量波取
\[
\delta\rho=0,\qquad \delta\bm{m}=0,\qquad
\delta(\rho\phi_a)=1 .
\]
\(\delta p=0\) 给出
\[
\Gammaone\,\delta(\rho E)+c^2\alpha_a=0,
\]
因此该波必须同时带有
\[
\delta(\rho E)=\chi_a=-\frac{\alpha_a c^2}{\Gammaone}.
\]
右特征向量为
\[
r_{\phi_a}=
\begin{bmatrix}
0\\0\\0\\0\\\chi_a\\0\\\vdots\\1\text{ at scalar }a\\\vdots\\0
\end{bmatrix}.
\]
代入压力微分可验证
\[
\delta p
=\Gammaone\chi_a+c^2\alpha_a=0.
\]
物种波的能量补偿为
\[
\boxed{\chi_s
=h_{\tr,s}-\frac{\gamma_{\tr}R_sT}{\Gammaone}},
\]
双温标量波的能量补偿为
\[
\boxed{\chi_{\ve}=1}.
\]
所以双温右特征向量为
\[
r_{E_{\ve}}=[0,0,0,0,1,0,\ldots,1\text{ at }\rho E_{\ve},\ldots,0]^T .
\]
对 \(a=i_{\ve}\)，\(\chi_{\ve}=1\) 且
\(\alpha_{\ve}=-\Gammaone/c^2\)。该标量波同时增加
\(\rho E\) 与 \(\rho E_{\ve}\)，使
\(\rho e_{\tr}=\rho E-\rho K-\rho E_{\ve}\) 和压力保持不变。

\section{特征值}

单位法向下的特征值已经由上节得到：
\[
\lambda_n^- = U-c,\qquad
\lambda_n^0=U,\qquad
\lambda_n^+ = U+c .
\]
\(\lambda_n^0\) 的重数为 \(M+3\)：两个切向动量波、一个接触波和
\(M\) 个 transported-scalar 波。曲线坐标方向满足
\(\FF_k=|\kk|\FF_n\)，所以
\[
\boxed{\lambda_- = U_k-c|\kk|,\qquad
\lambda_0=U_k,\qquad
\lambda_+=U_k+c|\kk|}.
\]

\section{左特征向量与投影}

对任意守恒扰动或通量型向量
\[
\zz=[z_\rho,z_{\rho u},z_{\rho v},z_{\rho w},z_{\rho E},
z_{\rho\phi_1},\ldots,z_{\rho\phi_M}]^T,
\]
并记 \(\zz_m=(z_{\rho u},z_{\rho v},z_{\rho w})^T\)。
定义
\[
\beta(\zz)=\sum_{a=1}^{M}\alpha_a z_{\rho\phi_a}.
\]
本节中 \(U=U_n=\nn\cdot\uu\)。
把 \(\zz\) 展开为
\[
\zz=a_-r_-+a_t r_t+a_b r_b+a_c r_c+\sum_a a_{\phi_a}r_{\phi_a}.
\]
由密度分量有
\[
z_\rho=a_-+a_c+a_+ .
\]
由标量分量有
\[
z_{\rho\phi_a}=\phi_a(a_-+a_c+a_+)+a_{\phi_a},
\]
因此
\[
\boxed{a_{\phi_a}=z_{\rho\phi_a}-\phi_a z_\rho}.
\]
切向动量直接给出
\[
\boxed{a_t=\bm{t}\cdot\zz_m-U_tz_\rho},\qquad
\boxed{a_b=\bm{b}\cdot\zz_m-U_bz_\rho}.
\]
压力微分作用在 \(\zz\) 上为
\[
\delta p(\zz)
=\Gammaone(z_{\rho E}-\uu\cdot\zz_m+Kz_\rho)
+c^2\beta(\zz).
\]
因为只有两个声波带有压力扰动，且
\(\delta p(r_-)=\delta p(r_+)=c^2\)，所以
\[
a_-+a_+=\frac{\delta p(\zz)}{c^2}.
\]
再由法向动量分量
\[
\nn\cdot\zz_m=Uz_\rho+c(a_+-a_-)
\]
得到
\[
a_+-a_-=\frac{\nn\cdot\zz_m-Uz_\rho}{c}.
\]
联立可得声波投影、接触投影和标量投影。写成行向量形式即
\[
L_-\zz=
\frac{1}{2}\left[
\frac{\Gammaone K+cU}{c^2}z_\rho
-\frac{(\Gammaone\uu+c\nn)\cdot\zz_m}{c^2}
+\frac{\Gammaone}{c^2}z_{\rho E}
+\beta(\zz)\right],
\]
\[
L_t\zz=-U_tz_\rho+\bm{t}\cdot\zz_m,\qquad
L_b\zz=-U_bz_\rho+\bm{b}\cdot\zz_m,
\]
\[
L_c\zz=
\left(1-\frac{\Gammaone K}{c^2}\right)z_\rho
+\frac{\Gammaone}{c^2}\uu\cdot\zz_m
-\frac{\Gammaone}{c^2}z_{\rho E}
-\beta(\zz),
\]
\[
L_+\zz=
\frac{1}{2}\left[
\frac{\Gammaone K-cU}{c^2}z_\rho
-\frac{(\Gammaone\uu-c\nn)\cdot\zz_m}{c^2}
+\frac{\Gammaone}{c^2}z_{\rho E}
+\beta(\zz)\right],
\]
\[
L_{\phi_a}\zz=z_{\rho\phi_a}-\phi_a z_\rho.
\]
这些行向量满足 \(\LL\RR=I\)。
双温标量系数 \(\alpha_{\ve}=-\Gammaone/c^2\) 代入后，
当 \(\zz=r_{E_{\ve}}\) 时
\[
L_-\zz=L_t\zz=L_b\zz=L_c\zz=L_+\zz=0,\qquad
L_{E_{\ve}}\zz=1.
\]
因此把 \(E_{\ve}\) 作为单独标量处理时，特征投影是严格闭合的。

\section{特征 WENO 离散}

在界面 \(i+1/2\) 处取 Roe 平均状态，构造局部特征矩阵
\(\LL_{i+1/2}\)、\(\RR_{i+1/2}\) 和最大谱半径
\[
\lambda_{\max}=\max_m\left(|U_k|+c|\kk|\right)_m .
\]
有限差分 WENO 使用局部 Lax--Friedrichs 分裂
\[
\FF^{\pm}_m=\frac{1}{2}\left(\FF_m\pm\lambda_{\max}\QQ_m\right),
\]
投影到特征空间
\[
\hat{\FF}^{\pm}_m=\LL_{i+1/2}\FF^{\pm}_m,
\]
分别做 WENO5 或 WENO7 重构：
\[
\hat{\FF}_{i+1/2}
=\mathcal{WENO}(\hat{\FF}^{+}_{i-r},\ldots,\hat{\FF}^{+}_{i+s})
+\mathcal{WENO}(\hat{\FF}^{-}_{i+1-r},\ldots,\hat{\FF}^{-}_{i+1+s}),
\]
最后回投影
\[
\FF_{i+1/2}=\RR_{i+1/2}\hat{\FF}_{i+1/2}.
\]
残差按有限差分形式写为
\[
\frac{\dd \QQ_i}{\dd t}
=-\left(\FF_{i+1/2}-\FF_{i-1/2}\right)+\cdots .
\]

\subsection{WENO 源码变量对应}

\texttt{src/WENO.cu} 的 x、y、z 三个 characteristic kernel 使用相同的局部特征结构。
共享数组 \texttt{Q} 存储守恒量，\texttt{F[0]} 存储 \(U_k\)，
\texttt{F[1:3]} 存储动量无粘通量，\texttt{F[4]} 存储压力，
\texttt{F[5]} 存储 \(c|\kk|\)，\texttt{F[6]} 存储 Jacobian。
\texttt{Q[5+i\_eve]} 是 \(\rho E_{\ve}\)，没有对 \(T_{\ve}\) 做特征重构；
\(T_{\ve}\) 只在源项和粘性项中通过辅助数组使用。界面 Roe 平均为
\[
\bar{\phi}_a=\texttt{svm[a]},\qquad
\bar{T}=\texttt{temp3},\qquad
\bar{c}=\texttt{cm},\qquad
\bar{\gamma}-1=\texttt{gm1}.
\]
\code{compute\_mixture\_characteristic\_thermo} 给出
\[
\texttt{scalar\_alpha[a]}=\alpha_a,\qquad
\texttt{energy\_coeff[a]}=\chi_a=-\frac{\alpha_a c^2}{\Gammaone}.
\]
对物种，\(\chi_s=h_{\tr,s}-\gamma_{\tr}R_sT/\Gammaone\)；对双温标量，
\(\chi_{\ve}=1\)。

声波和接触波投影中的标量压力贡献由
\[
\beta(\QQ)=\sum_a\alpha_a Q_{5+a},\qquad
\beta(\FF)=\sum_a\alpha_a Q_{5+a}U_k
\]
给出。物理通量中的标量分量满足
\(F_{5+a}=Q_{5+a}U_k\)，因此 \(\beta(\FF)\) 正是
\(\sum_a\alpha_aF_{5+a}\)。代码中对应为
\[
\code{sumBetaQ[m]}=\sum_a\code{scalar\_alpha[a]}\,\code{Q[5+a][m]},
\]
\[
\code{sumBetaF[m]}=\sum_a\code{scalar\_alpha[a]}\,\code{Q[5+a][m]}\,\code{F[0][m]}.
\]
五个气动特征行使用 \texttt{L[0:4]} 表示守恒量
\((\rho,\rho u,\rho v,\rho w,\rho E)\) 的系数，并通过
\texttt{temp1} 加入标量部分：声波 \texttt{temp1=0.5}，接触波
\texttt{temp1=-1}，两个切向波 \texttt{temp1=0}。这与左特征向量中的
\[
L_-\sim +\frac{1}{2}\beta,\qquad
L_c\sim -\beta,\qquad
L_+\sim +\frac{1}{2}\beta
\]
一致。每个标量波的投影单独计算为
\[
L_{\phi_a}\QQ=Q_{5+a}-\bar{\phi}_aQ_0,
\qquad
L_{\phi_a}\FF=U_k(Q_{5+a}-\bar{\phi}_aQ_0),
\]
对应代码中的
\code{Q[5+l]-svm[l]*Q[0]} 和
\code{F[0]*(Q[5+l]-svm[l]*Q[0])}。

回投影时，代码先形成
\[
\texttt{temp1}=f_-+f_c+f_+,\qquad
\texttt{temp3}=f_- - f_+,
\]
并写出质量、动量和不含独立标量波补偿的能量贡献。该能量基础项对应
\[
\bar H(f_-+f_c+f_+)-\bar U\bar c(f_- - f_+)
 +U_t f_t+U_b f_b-\frac{\bar c^2}{\bar\Gamma}f_c .
\]
随后标量回投影为
\[
F_{5+a}=\bar{\phi}_a\,\texttt{temp1}+f_{\phi_a},
\]
对应 \texttt{fc/gc/hc(...,5+l)=svm[l]*temp1+fChar[5+l]}；总能量再加
\[
\sum_a \chi_a f_{\phi_a},
\]
对应 \code{sum energy\_coeff[l]*fChar[5+l]}。当 \(a=i_{\ve}\) 时
\(\chi_{\ve}=1\)，所以 \(\rho E_{\ve}\) 标量波的能量补偿被等量加入总能量。
Roe 矩阵耗散路径 \texttt{src/InviscidScheme.cu} 使用同一组
\texttt{scalar\_alpha} 和 \texttt{energy\_coeff}：\code{c1} 中累加
\(\sum_a\alpha_a\Delta q_a\)，回投影时 \code{c3} 累加
\(\sum_a\chi_a b_{5+a}\)。因此 WENO characteristic kernel 与 Roe
耗散路径对 \(E_{\ve}\) 标量的处理一致。

\section{源码映射与实现核对}

\begin{longtable}{>{\raggedright\arraybackslash}p{0.27\linewidth}
                  >{\raggedright\arraybackslash}p{0.25\linewidth}
                  >{\raggedright\arraybackslash}p{0.40\linewidth}}
\hline
数学项 & 源码位置 & 代码实现对应\\
\hline
\endhead
编译期双温选择 &
\code{CMakeLists.txt}, \code{src/Define.h} &
\code{COREFL\_ENABLE\_TWO\_TEMPERATURE} 生成
\texttt{kTwoTemperature}；设备端双温路径均在
\texttt{if constexpr (kTwoTemperature)} 下编译。\\
\hline
\(\rho E_{\ve}\) 布局 &
\code{src/Parameter.cpp} &
\code{i\_eve=n\_spec}, \code{i\_eve\_cv=5+n\_spec}；
\texttt{n\_scalar\_transported} 加 1。\texttt{Tve} 只作为
\texttt{n\_other\_var} 输出和 \texttt{temperature\_ve} 辅助场。\\
\hline
总能量闭合 \(E=K+\sum_sY_se_{\tr,s}(T)+E_{\ve}\) &
\code{src/FieldOperation.cuh}, \code{src/FieldOperation.cu}, \code{src/Thermo.cu} &
\code{compute\_total\_energy} 在 NASA 单温内能
\(h_{eq,s}-R_sT\) 后减 \(\Mve(T,Y)\)，再加
\code{sv[i\_eve]}。\code{compute\_mixture\_tr\_energy} 返回
\(e_{\tr}(T,Y)\)，反算温度时固定 \texttt{Eve}。\\
\hline
无粘通量 \(\rho E_{\ve}U_k\) &
\code{src/WENO.cu}, \code{src/InviscidScheme.cu} &
\texttt{WENO.cu} 将 \code{Q[5+l]} 作为所有 transported scalar 的守恒量，
标量特征通量使用 \code{F[0]*(Q[5+l]-svm[l]*Q[0])}；
物理通量层面即 \(Q_{5+a}U_k\)。当 \(a=i_{\ve}\) 时即
\(\rho E_{\ve}U_k\)。\\
\hline
\(\alpha_s,\alpha_{\ve}\) 和 \(\chi_a\) &
\code{src/Thermo.cuh} &
\code{compute\_mixture\_characteristic\_thermo} 中
\code{h\_tr\_i=h\_eq-eve\_eq},
\code{alpha=(gamma*R\_s*T-gm1*h\_tr\_i)/c2}；
\code{scalar\_alpha[i\_eve]=-gm1/c2}；
\code{energy\_coeff[l]=-alpha*c2/gm1}，故
\code{energy\_coeff[i\_eve]=1}。\\
\hline
WENO 特征投影/回投影 &
\code{src/WENO.cu} &
x/y/z 三个 characteristic kernel 均用左右界面状态构造 Roe 平均
\texttt{svm} 和 \texttt{temp3}，调用
\code{compute\_mixture\_characteristic\_thermo}。投影时
\code{sumBetaQ=sum scalar\_alpha[n]*Q[5+n]}，
\code{sumBetaF=sum scalar\_alpha[n]*Q[5+n]*F[0]}；
回投影时能量分量加
\code{sum energy\_coeff[l]*fChar[5+l]}，
标量分量为 \code{svm[l]*temp1+fChar[5+l]}。\\
\hline
化学源 \(\sum_s\dot\omega_s \varepsilon_{\ve,s}(T_{\ve})\) &
\code{src/FiniteRateChem.cu} &
\code{compute\_two\_temperature\_chemical\_source} 读取
\code{zone->temperature\_ve}，循环
\code{omega[l]*compute\_ve\_energy(l,tve,param)}；
有限率化学 kernel 将结果加到 \code{dq(...,5+i\_eve)}。
瞬态 RK3 在 \texttt{Parameter.cpp} 中强制 \code{chemSrcMethod=0}，
因此 \(\rho E_{\ve}\) 化学源在当前双温 CFD 路径为显式。\\
\hline
VT 源 \(\dot\Theta^{\tr:\ve}\) &
\code{src/Thermo.cu}, \code{src/FieldOperation.cuh}, \code{src/RK.cuh}, \code{src/ConstVolumeReactor.cpp} &
\code{compute\_vt\_relaxation\_source} 对
\code{theta\_v[i]>0} 的物种累加
\(\rho Y_i[\varepsilon_{\ve,i}(T)-\varepsilon_{\ve,i}(T_{\ve})]/\tau_i\)，
其中两端均调用 \code{compute\_ve\_energy}。
\code{apply\_vt\_relaxation\_1\_point} 用
\code{exp(-dt/tau\_eff)} 更新 \texttt{rhoEve}，再反解
\texttt{Tve} 并重新闭合 \(T,p\)。\\
\hline
粘性热通量 &
\code{src/FieldOperation.cuh}, \code{src/ViscousScheme.cu} &
\code{compute\_transport\_property} 后保存的平衡导热为
\code{conductivity\_eq}；随后
\code{thermal\_conductivity=conductivity\_eq*cp\_tr/cp\_mix\_tr}，
\code{thermal\_conductivity\_ve=conductivity\_eq*cv\_ve/cp\_mix\_tr}。
8 阶 collocated 通量在总能量 \texttt{f/g/hFlux(...,4)} 中加入
\(\kappa_{\ve}\nabla T_{\ve}\)，并在
\code{f/g/hFlux(...,5+i\_eve)} 写入 \(\rho E_{\ve}\) 粘性通量。\\
\hline
物种扩散携带能量 &
\code{src/ViscousScheme.cu}, \code{src/Thermo.cuh} &
总能量物种扩散焓由
\code{compute\_nonequilibrium\_diffusion\_enthalpy} 给出
\(h_{eq,s}(T)-\varepsilon_{\ve,s}(T)+\varepsilon_{\ve,s}(T_{\ve})\)。
\(\rho E_{\ve}\) 通量循环中使用
\code{compute\_ve\_energy(l,tve\_m)} 乘同一物种扩散通量。\\
\hline
wall 热流输出 &
\code{src/PostProcess.cu} &
\texttt{q\_tr\_face} 使用 \(\kappa_{\tr}\nabla T\) 和
\(h_{eq,s}(T)-\varepsilon_{\ve,s}(T)\)；\texttt{q\_ve\_face} 使用
\(\kappa_{\ve}\nabla T_{\ve}\) 和
\(\varepsilon_{\ve,s}(T_{\ve})\)。输出
\texttt{q\_tr}, \texttt{q\_ve}, \texttt{q\_total=q\_tr+q\_ve}。\\
\hline
restart/output &
\code{src/Initialize.cu}, \code{src/FieldIO.h}, \code{src/Field.cu} &
\texttt{Parameter.cpp} 将 \texttt{Eve} 加入 scalar 变量名，将
\texttt{Tve} 加入额外输出变量名；restart/边界同步使用
\(\Mve(T_{\ve},Y)=E_{\ve}\)。\\
\hline
\end{longtable}

\section{VT 松弛更新}

令参与 Landau--Teller/Park 松弛的物种集合为
\[
\mathcal{V}=\{s:\theta_{v,s}>0\}.
\]
COREFL 为这些有分子振动特征温度的物种构造 VT 松弛时间。
对 \(s\in\mathcal{V}\)，松弛源采用完整的
\(\varepsilon_{\ve,s}=\varepsilon_{vib,s}+\varepsilon_{el,s}\) 差值：
\[
\dot\Theta^{\tr:\ve}
=\sum_{s\in\mathcal{V}}\rho Y_s
\frac{\varepsilon_{\ve,s}(T)-\varepsilon_{\ve,s}(T_{\ve})}{\tau_s}.
\]
代码先把这个多物种源等效成
\[
\tau_{eff}
=\frac{\rho\left(\Mve^{eq,\mathcal{V}}-\Mve^{\mathcal{V}}\right)}
{\dot\Theta^{\tr:\ve}},
\]
其中
\[
\Mve^{eq,\mathcal{V}}=\sum_{s\in\mathcal{V}}Y_s\varepsilon_{\ve,s}(T),
\qquad
\Mve^{\mathcal{V}}=\sum_{s\in\mathcal{V}}Y_s\varepsilon_{\ve,s}(T_{\ve}).
\]
然后在每个 RK stage 后对参与松弛的 VE 能量做指数更新
\[
\mathcal{M}_{\ve,new}^{\mathcal{V}}
=\Mve^{\mathcal{V}}\exp(-\Delta t/\tau_{eff})
+\Mve^{eq,\mathcal{V}}\left[1-\exp(-\Delta t/\tau_{eff})\right].
\]
守恒量更新为
\[
(\rho E_{\ve})_{new}
=(\rho E_{\ve})_{old}
+\rho\left(\mathcal{M}_{\ve,new}^{\mathcal{V}}-\Mve^{\mathcal{V}}\right).
\]
随后代码由新的混合物 \(E_{\ve}\) 反解 \(T_{\ve}\)，并在总能量不变的条件下重新闭合
\(T,p\)。

\(\mathcal{V}\) 之外的原子电子能级没有分子振动松弛时间，
因此不进入 LT/Park 松弛循环；但它们仍然通过
\(\varepsilon_{\ve,s}(T_{\ve})\) 进入守恒量、化学源、粘性扩散和 \(T_{\ve}\) 反解，
也通过 \(\varepsilon_{\ve,s}(T)\) 进入 NASA 热力学拆分。

\section{振动--电子能量使用点核对}

为避免 \(\varepsilon_{\ve,s}(T)\) 和 \(\varepsilon_{\ve,s}(T_{\ve})\) 混淆，所有 VE 相关位置按下表核对。

\begin{longtable}{>{\raggedright\arraybackslash}p{0.29\linewidth}
                  >{\raggedright\arraybackslash}p{0.31\linewidth}
                  >{\raggedright\arraybackslash}p{0.32\linewidth}}
\hline
用途 & 正确数学形式 & 源码结论\\
\hline
\endhead
从 NASA 单温热力学定义 \(e_{\tr}\) &
\(e_{\tr,s}(T)=h_{eq,s}(T)-R_sT-\varepsilon_{\ve,s}(T)\) &
\code{compute\_mixture\_tr\_energy} 和特征热力学均在 \(T\) 下减去 \(\Mve(T,Y)\)。\\
\hline
守恒变量 &
\(\rho E_{\ve}\)，其中 \(E_{\ve}\) 是独立 transported scalar &
\code{sync\_two\_temperature\_state\_from\_sv} 和 restart/边界通过
\(\Mve(T_{\ve},Y)=E_{\ve}\) 保持 \texttt{Eve} 与 \(T_{\ve}\) 一致。\\
\hline
总能量 &
\(E=K+\sum_sY_se_{\tr,s}(T)+E_{\ve}\) &
\code{compute\_total\_energy} 先减 \(\Mve(T,Y)\)，再加 \code{sv[i\_eve]}；\code{sv[i\_eve]} 是独立 \(E_{\ve}\)。\\
\hline
由 \(\rho E\) 反算 \(T\) &
\(e_{\tr}(T,Y)=E-K-E_{\ve}\) &
\code{compute\_temperature\_and\_pressure} 固定 \texttt{Eve}，
用 \code{compute\_mixture\_tr\_energy(T)} 反解 \(T\)。\\
\hline
无粘通量 &
\(\rho E_{\ve}U_k\) &
标量通量统一为 \(Q_{5+a}U_k\)，\texttt{i\_eve} 对应 \(\rho E_{\ve}\)。\\
\hline
化学源 &
\(\sum_s\dot\omega_s\varepsilon_{\ve,s}(T_{\ve})\) &
\texttt{FiniteRateChem.cu} 和 host reactor 化学源均用 \texttt{compute\_ve\_energy(l,tve)}。\\
\hline
VT 松弛源 &
\(\sum_{s\in\mathcal{V}}\rho Y_s[\varepsilon_{\ve,s}(T)-\varepsilon_{\ve,s}(T_{\ve})]/\tau_s\) &
主 CFD 和 host reactor 均只对 \(\theta_{v,s}>0\) 的参与物种做 LT 松弛，但差值使用完整 \code{compute\_ve\_energy}，不再只用振动能。\\
\hline
\(\rho E_{\ve}\) 粘性扩散 &
\(\kappa_{\ve}\nabla T_{\ve}+\sum_sD_s^{code}\varepsilon_{\ve,s}(T_{\ve})\) &
\texttt{ViscousScheme.cu} 三个方向均用 \(T_{\ve}\) 面值计算 \code{compute\_ve\_energy(l,tve\_m)}。\\
\hline
总能量物种扩散 &
\(\sum_sD_s^{code}[h_{\tr,s}(T)+\varepsilon_{\ve,s}(T_{\ve})]\) &
\code{compute\_nonequilibrium\_diffusion\_enthalpy} 实现为 \(h_{eq,s}(T)-\varepsilon_{\ve,s}(T)+\varepsilon_{\ve,s}(T_{\ve})\)。\\
\hline
wall 热流拆分 &
\(q_{\tr}\) 使用 \(h_{\tr,s}(T)\)，\(q_{\ve}\) 使用 \(\varepsilon_{\ve,s}(T_{\ve})\) &
\texttt{PostProcess.cu} 中 \(q_{\tr}\) 的物种扩散焓减 \(\varepsilon_{\ve,s}(T)\)，\(q_{\ve}\) 的物种扩散项加 \(\varepsilon_{\ve,s}(T_{\ve})\)。\\
\hline
\end{longtable}

\section{粘性项符号核对}

SU2 文档以物理物种扩散通量 \(\bm{J}_s\) 写 Fick 定律。
COREFL 的物种方程残差把存储的粘性通量散度直接加到 \(\dd\QQ/\dd t\)，
因此代码中
\[
\bm{D}_s^{code}=-\bm{J}_s .
\]
在该符号约定下，双温方程的粘性通量应为
\[
\bm{F}_{\rho E_{\ve}}^{vis}
=\kappa_{\ve}\nabla T_{\ve}
+\sum_s \bm{D}_s^{code}\varepsilon_{\ve,s}(T_{\ve}),
\]
总能量粘性通量应为
\[
\bm{F}_{\rho E}^{vis}
=\bm{\tau}\cdot\uu+\kappa_{\tr}\nabla T+\kappa_{\ve}\nabla T_{\ve}
+\sum_s \bm{D}_s^{code}
\left[h_{\tr,s}(T)+\varepsilon_{\ve,s}(T_{\ve})\right].
\]
代码采用如下导热拆分：
\[
\kappa_{\tr}=\kappa_{eq}\frac{c_{p,\tr}}{c_{p,eq}},
\qquad
\kappa_{\ve}=\kappa_{eq}\frac{c_{v,\ve}(T_{\ve})}{c_{p,eq}},
\]
其中 \(c_{p,\tr}=c_{p,eq}-c_{v,\ve}(T)\)。源码中的
\code{cp\_mix\_tr} 在拆分前是平衡混合物 \(c_{p,eq}\)，
\code{cp\_tr} 是减去平衡 VE 热容后的平动--转动部分。
这与 SU2 文档中将不同能量模态导热分开处理的要求一致。

源码中 8 阶 collocated 粘性通量和 2 阶面通量的数组索引不同：
\texttt{fFlux/gFlux/hFlux} 按守恒变量索引写入，因此总能量为
\texttt{(...,4)}，\(\rho E_{\ve}\) 为 \texttt{(...,5+i\_eve)}；
\texttt{vis\_flux} 只存动量、能量和标量粘性通量，能量为
\texttt{(...,3)}，第 \(l\) 个标量为 \texttt{(...,4+l)}。
对应代码为
\[
\begin{array}{c|c|c}
\text{离散路径} & \rho E\text{ 粘性通量} & \rho E_{\ve}\text{ 粘性通量}\\ \hline
8\text{ 阶 collocated}
& \texttt{f/g/hFlux(...,4)}
& \code{f/g/hFlux(...,5+i\_eve)}\\
2\text{ 阶 face}
& \texttt{vis\_flux(...,3)}
& \texttt{vis\_flux(...,4+i\_eve)}
\end{array}
\]
两条路径都先把 \(\kappa_{\ve}\nabla T_{\ve}\) 加入总能量粘性通量，
再把同一导热项作为 \(\rho E_{\ve}\) 方程的导热部分。总能量方程表示
所有能量模态之和，\(\rho E_{\ve}\) 方程表示其中一个独立模态；
同一物理 VE 导热通量同时出现在两个方程中，不构成总能量内的重复相加。
物种扩散部分分别使用
\(h_{\tr,s}(T)+\varepsilon_{\ve,s}(T_{\ve})\) 和
\(\varepsilon_{\ve,s}(T_{\ve})\)。因此总能量包含实际扩散携带的
VE 能量一次，\(\rho E_{\ve}\) 方程也得到独立的 VE 扩散通量。

\section{运行约束}

双温路径限定为瞬态 RK3：
\[
\texttt{steady=0},\qquad \texttt{temporal\_scheme=3},\qquad
\texttt{species=1},\qquad \texttt{reaction=1}.
\]
该限制在 \texttt{src/Parameter.cpp} 中强制执行。
在该路径下 \texttt{Parameter.cpp} 同时把 \code{chemSrcMethod} 置为 0；
\(\rho E_{\ve}\) 化学源随 RK3 残差显式更新，VT 松弛源在每个 RK stage
后通过指数公式单独更新。
默认热力学和反应模型保持 NASA9 与 \texttt{Combustion2Part}。

\end{document}
